\clearpage
\begin{summarybox}
\textbf{\textcolor{CSummary}{这条结论为什么存在？}}

因为计算 $f(g(x))$ 要\textbf{先经过 $g$，再经过 $f$}。
与其硬看复杂式子，不如逐层判断大小关系怎么变化。

\medskip
\textbf{\textcolor{CSummary}{一句话记住}}

\textbf{递增保留大小顺序，递减反转大小顺序；两次反转仍增，一次反转就减。}
即“同增异减”：增增得增，减减得增，增减或减增得减。

\medskip
\textbf{\textcolor{CSummary}{四步用法}}
\begin{enumerate}[leftmargin=*, itemsep=0.25em]
\item 先看整个函数在什么 $x$ 上有定义。
\item 设 $t=g(x)$，把里外分清。
\item 把 $x$ 分成若干段：每段里层保持同一方向，外层对实际得到的 $t$ 也保持同一方向。
\item 逐段套口诀，把答案写回 $x$ 的区间。
\end{enumerate}

\medskip
\textbf{\textcolor{CSummary}{三个能举一反三的发现}}
\begin{itemize}[leftmargin=*, itemsep=0.25em]
\item \textbf{不求导也能判增减：}比较任意两个 $x_1<x_2$ 的函数值就够了。
\item \textbf{外层若递减：}内层取到最小值的位置，整体反而取到最大值。
\item \textbf{多层也通用：}数递减层数，偶数次反转得增，奇数次反转得减。
\end{itemize}

\medskip
\textbf{\textcolor{CSummary}{最后的底线}}

\textbf{只查定义域不够。} 还要确认外层在内层\emph{实际输出的范围内}单调。
如果某层只是“不减或不增”，最后不能擅自说“严格递增或严格递减”。
\end{summarybox}
